\chapter{Planificación} Para gestionar la planificación del proyecto, se ha utilizado el método Kanban para la organización y asignación de tareas, y la herramienta Wakatime para registrar y analizar el tiempo invertido en cada actividad. El objetivo es tener un seguimiento preciso de las tareas completadas y del esfuerzo dedicado, asegurando el cumplimiento de los objetivos semanales en un plazo de seis semanas.

\section{Planificación semanal}
Dado el tiempo total de un mes y medio para la realización del proyecto, la planificación se ha organizado de la siguiente manera:

\begin{itemize}
    \item \textbf{Primera semana}: documentación y análisis del estado del proyecto.
    \item \textbf{Segunda semana}: resolución de deuda técnica.
    \item \textbf{Tercera semana}: desarrollo de la interfaz de usuario.
    \item \textbf{Cuarta semana}: redacción de la memoria.
    \item \textbf{Quinta semana}: incorporación de funcionalidades adicionales y tests al proyecto. Documentación adicional.
    \item \textbf{Sexta semana}: solución de problemas de seguridad y accesibilidad, así como finalización de documentación y ajustes finales.
\end{itemize}
Cada sección semanal incluirá:
\begin{itemize}
    \item Lista de tareas (Kanban)
    \item Registro de tiempo (Wakatime)
    \item Literatura y referencias empleadas
    \item Commits de github (si aplica)
\end{itemize}


\section{Semana 1: Documentación y prueba del estado del proyecto}
En la primera semana, se dedicará tiempo a la comprensión profunda del proyecto, evaluando su estado actual y revisando la documentación existente para identificar posibles problemas técnicos y áreas de mejora.

\begin{itemize}
    \item Tareas de la semana:
          \begin{itemize}
              \item Revisión de la estructura del proyecto y código existente.
              \item Documentación del estado actual del proyecto.
              \item Identificación de áreas con deuda técnica para su futura resolución.
              \item Configuración de herramientas de seguimiento (Kanban y Wakatime). \item Revisión de literatura sobre MLOps y buenas prácticas de desarrollo de software.
          \end{itemize}
    \item Kanban semanal: Imagen con el tablero Kanban correspondiente a las tareas de la semana.
    \item Wakatime semanal: Imagen del resumen semanal de tiempo trabajado en Wakatime.
    \item Literatura utilizada:
          \begin{itemize}
              \item \cite{mlops} MLOps: buenas prácticas y fundamentos.
              \item \cite{introtodebt} Introducción al concepto de deuda técnica.
              \item \cite{DORA} Informe DORA sobre rendimiento de equipos de desarrollo.
          \end{itemize}
\end{itemize}

\section{Semana 2: Resolución de deuda técnica}
Durante la segunda semana, el enfoque será identificar y corregir la deuda técnica en el proyecto, aplicando buenas prácticas de programación para mejorar la calidad y mantenibilidad del código.

\begin{itemize}
    \item Tareas de la semana:
          \begin{itemize}
              \item Refactorización de módulos con código duplicado o complejo.
              \item Optimización de funciones para mejorar el rendimiento.
              \item Documentación de los cambios realizados y su justificación.
          \end{itemize}
    \item Kanban semanal: Imagen del tablero Kanban de la semana.
    \item Wakatime semanal: Imagen del resumen de tiempo en Wakatime.
    \item Literatura utilizada:
          \begin{itemize}
              \item \cite{cunningham} Artículo de Ward Cunningham sobre deuda técnica. \item \cite{introtodebt} Conceptos y gestión de deuda técnica. \end{itemize}
\end{itemize}

\section{Semana 3: Desarrollo de la interfaz de usuario}
En la tercera semana, el objetivo será desarrollar una interfaz de usuario intuitiva para el sistema, basada en las prácticas de usabilidad y accesibilidad.

\begin{itemize}
    \item Tareas de la semana:
          \begin{itemize}
              \item Diseño de los elementos de la interfaz (formularios, botones, visualización de logs).
              \item Implementación de la interfaz utilizando React y Bootstrap.
              \item Pruebas de usabilidad y accesibilidad en la interfaz desarrollada. \item Documentación de la estructura y funcionalidades de la interfaz.
          \end{itemize}
    \item Kanban semanal: Imagen del tablero Kanban de la semana.
    \item Wakatime semanal: Imagen del resumen de tiempo en Wakatime.
    \item Literatura utilizada:
          \begin{itemize}
              \item \cite{reactdocs} Documentación de React y Bootstrap.
              \item \cite{accessibility} Guías de accesibilidad y usabilidad para interfaces web.
          \end{itemize}
\end{itemize}

\section{Semana 4: Redacción de la memoria}
En la cuarta semana, el objetivo principal será avanzar en la memoria del proyecto, documentando todos los aspectos técnicos y funcionales, así como el proceso de desarrollo.

\begin{itemize}
    \item Tareas de la semana:
          \begin{itemize}
              \item Redacción de la sección de introducción y objetivos del proyecto. \item Documentación de los requisitos y la especificación del sistema. \item Descripción detallada de las soluciones aplicadas para resolver la deuda técnica.
              \item Redacción de la planificación semanal y metodología.
          \end{itemize}
    \item Kanban semanal: Imagen del tablero Kanban de la semana.
    \item Wakatime semanal: Imagen del resumen de tiempo en Wakatime.
    \item Literatura utilizada:
          \begin{itemize}
              \item Documentación interna del proyecto.
              \item Cualquier fuente citada en el TFG, para una referencia precisa en el documento.
          \end{itemize}
\end{itemize}

\section{Semana 5: Funcionalidades adicionales y pruebas}
En la quinta semana se añadirán funcionalidades adicionales y se implementarán pruebas para mejorar la calidad y estabilidad del sistema, además de completar documentación técnica.

\begin{itemize}
    \item Tareas de la semana:
          \begin{itemize}
              \item Implementación de funcionalidades adicionales solicitadas por el tutor.
              \item Creación de pruebas unitarias y de integración para el nuevo código.
              \item Actualización de la documentación con las nuevas funcionalidades y pruebas.
          \end{itemize}
    \item Kanban semanal: Imagen del tablero Kanban de la semana.
    \item Wakatime semanal: Imagen del resumen de tiempo en Wakatime.
    \item Literatura utilizada:
          \begin{itemize}
              \item Referencias adicionales sobre MLOps y diseño de pruebas en proyectos de software.
          \end{itemize}
\end{itemize}

\section{Semana 6: Solución de problemas de seguridad y accesibilidad}
En la última semana, se centrará en resolver problemas de seguridad y accesibilidad en el sistema, además de completar la documentación final del proyecto.

\begin{itemize}
    \item Tareas de la semana:
          \begin{itemize}
              \item Identificación y corrección de posibles vulnerabilidades de seguridad en la interfaz.
              \item Validación de la accesibilidad de la interfaz para usuarios con diferentes necesidades.
              \item Revisión final de toda la documentación del proyecto.
          \end{itemize}
    \item Kanban semanal: Imagen del tablero Kanban de la semana.
    \item Wakatime semanal: Imagen del resumen de tiempo en Wakatime.
    \item Literatura utilizada:
          \begin{itemize}
              \item \cite{security} Guías sobre seguridad en desarrollo web.
              \item \cite{accessibility} Guías de accesibilidad y usabilidad para interfaces web.
          \end{itemize}
\end{itemize}